% Complete Section 4.2, original PDF 45-47 / printed 35-37.
% mathsf H is the non-Fraktur harmonic space; mathcal H remains cohomology.
% block: harmonic-hodge-poincare-heading page:45 type:section
\subsection{{Harmonic forms, the Hodge decomposition, and the Poincaré inequality}}\label{sec:harmonic-hodge-poincare}

% block: hilbert-structure-consequences page:45 type:ordinary-paragraph
{Next we derive several simple, but crucial, consequences of the Hilbert space structure. In this section, we assume that $(W,d)$ is a Hilbert complex, which is not necessarily closed.}

% block: harmonic-forms-heading page:45 type:subsection
\subsubsection{{Harmonic forms}}\label{subsec:harmonic-forms}

% block: harmonic-form-definition page:45 type:definition
\Needspace{6\baselineskip}
\begin{Definition}\label{def:harmonic-form}
{An element $v\in W^k$ is called \emph{harmonic} or a \emph{harmonic $k$-form} if both $dv=0$ and $d^*v=0$.}
\end{Definition}

% block: harmonic-space-definition page:45 type:ordinary-paragraph
{We denote the space of harmonic $k$-forms by $\mathsf H^k$, so $\mathsf H^k=\mathcal Z^k\cap\mathcal Z_k^*$. By the first assertion of Theorem~\ref{thm:adjoint-kernel-range-orthogonality}, we may as well write}
\begin{equation}\label{eq:harmonic-cocycle-orthogonality}
\mathsf H^k=\mathcal Z^k\cap\mathcal B^{k,\perp},
\end{equation}

% block: harmonic-reduced-cohomology page:45 type:ordinary-paragraph
{meaning that the harmonic forms are the cocycles which are orthogonal to coboundaries. If the coboundary space is not closed, we may close it and equivalently say that $\mathsf H^k$ is the orthogonal complement of the closed subspace $\overline{\mathcal B}^k$ inside the Hilbert space $\mathcal Z^k$. We may then canonically identify the quotient space, which is the reduced cohomology space $\mathcal Z^k/\overline{\mathcal B}^k$, with $\mathsf H^k$, which is a subspace of $W^k$. Of course, for a closed Hilbert complex, we do not need to specify ``reduced,'' and the result is simply}
\begin{equation}\label{eq:harmonic-cohomology-isomorphism}
\mathsf H^k\cong\mathcal H^k.
\end{equation}

% block: harmonic-canonical-representatives pages:45,46 type:ordinary-paragraph
{Thus the harmonic forms give a concrete realization of the cohomology of the complex as a closed subspace of the Hilbert space $W^k$. This is illustrated in Figure~\ref{fig:cohomology-harmonic-representatives}. The elements of the cohomology space $\mathcal H^k$ are the equivalence classes of elements of $\mathcal Z^k$ modulo $\mathcal B^k$. These provide a decomposition of $\mathcal Z^k$ into closed affine subspaces obtained from $\mathcal B^k$ by translation. Each such subspace contains a distinguished element, namely, the unique element orthogonal to $\mathcal B^k$, which is also characterized as the unique element of the smallest norm. These are the harmonic $k$-forms.}

% block: cohomology-geometry-figure page:45 type:figure
\begin{figure}[htbp]\centering
\input{figures/cohomology-geometry.tex}
\caption{{The space $\mathcal B^k$ of coboundaries is a subspace of the space $\mathcal Z^k$ of cocycles. The cohomology space consists of equivalence classes of cocycles, with the class $[z]$ of a cocycle $z$ being the affine subspace through $z$ parallel to $\mathcal B^k$. Each cohomology class $[z]$ is represented by a unique element $z_0$ belonging to $\mathsf H^k$, which is its element of its smallest norm.}}\label{fig:cohomology-harmonic-representatives}
\end{figure}

% block: compactness-property-intro page:46 type:ordinary-paragraph
{We may use harmonic forms to obtain a useful criterion for a Hilbert complex to be Fredholm, i.e., for its homology spaces to be finite. The space $V^k\cap V_k^*\subset W^k$ is a Hilbert space when endowed with the norm}
\[
\|v\|_{V\cap V^*}^2=\|v\|_V^2+\|v\|_{V^*}^2=2\|v\|^2+\|dv\|^2+\|d^*v\|^2.
\]

% block: hilbert-norm-conventions page:46 type:ordinary-paragraph
{Here and below, we write $\|\cdot\|_V$ for the norm in each of the spaces $V^k$, $\|\cdot\|_{V^*}$ for the norm in the spaces $V_k^*$, and simply $\|\cdot\|$ for the norm in the spaces $W^k$.}

% block: compactness-property-definition page:46 type:definition
\Needspace{6\baselineskip}
\begin{Definition}\label{def:hilbert-complex-compactness-property}
{A Hilbert complex is said to have the compactness property if the inclusion $V^k\cap V_k^*\subset W^k$ is compact for each $k$.}
\end{Definition}

% block: compactness-implies-fredholm page:46 type:theorem
\Needspace{6\baselineskip}
\begin{Theorem}\label{thm:compactness-implies-fredholm}
{If a Hilbert complex has the compactness property, then it is closed and Fredholm.}
\end{Theorem}

% block: compactness-closed-range-proof page:46 type:proof
\begin{Proof}
{We first check that $d$ has closed range. Let $T$ be the restriction of $d$ to $D(T):=V\cap\mathcal Z^\perp$, viewed as an unbounded operator $W\to W$. It is easy to check that the range of $T$ coincides with the range of $d$ and, since $d$ is a closed operator, $T$ is closed as well. Since $\mathcal Z^\perp=\overline{\mathcal B}^*\subset\mathcal Z^*\subset V^*$, $D(T)\subset V\cap V^*$ and the graph norm on $D(T)$ is equivalent to the $V\cap V^*$ norm. By Theorem~\ref{thm:compact-domain-closed-range}, $T$, and thus $d$, has closed range.}

% block: compactness-finite-dimensional-proof page:46 type:proof
{Next, we note that, on $\mathsf H\subset V\cap V^*$, the $V\cap V^*$ norm is equivalent to the $W$ norm. Thus the closed unit ball in $\mathsf H$ is compact, which implies that $\mathsf H\cong\mathcal H$ is finite dimensional, i.e., the complex is Fredholm.}
\end{Proof}

% block: hodge-decomposition-heading page:46 type:subsection
\subsubsection{{Hodge decomposition}}\label{subsec:hodge-decomposition}

% block: hodge-helmholtz-intro page:46 type:ordinary-paragraph
{The Hilbert complex structure leads directly to the Hodge decomposition, a generalization of the Helmholtz decomposition of vector calculus.}

% block: cocycle-orthogonal-decomposition page:46 type:ordinary-paragraph
{Since the space of harmonic $k$-forms is the orthogonal complement of $\overline{\mathcal B}^k$ within $\mathcal Z^k$, it follows immediately that}
\begin{equation}\label{eq:cocycle-harmonic-decomposition}
\mathcal Z^k=\overline{\mathcal B}^k\BookOrthSum\mathsf H^k,
\end{equation}

% block: orthogonal-sum-definition page:46 type:ordinary-paragraph
{where the symbol $\BookOrthSum$ indicates a sum which is orthogonal (note that in this context we need not distinguish between orthogonality in $W$ and orthogonality in $V$). Since $\mathcal Z^k$ is a closed subspace of both $W^k$ and $V^k$, we have}
\[
W^k=\mathcal Z^k\BookOrthSum\mathcal Z^{k\perp}=\mathcal Z^k\BookOrthSum\overline{\mathcal B}_k^*,\qquad V^k=\mathcal Z^k\BookOrthSum\mathcal Z^{k\perp V},
\]
{where $\mathcal Z^{k\perp V}=\mathcal Z^{k\perp}\cap V=\overline{\mathcal B}_k^*\cap V$. Combining with \eqref{eq:cocycle-harmonic-decomposition} we obtain the Hodge decomposition.}

% block: hodge-decomposition-theorem page:47 type:theorem
\Needspace{8\baselineskip}
\begin{Theorem}{Hodge {decomposition}}\label{thm:hodge-decomposition}
{Any Hilbert complex gives rise to the orthogonal decompositions}
\begin{equation}\label{eq:hodge-base-decomposition}
\begin{tikzpicture}[baseline=(m.base)]
\matrix (m) [matrix of math nodes,ampersand replacement=\&,column sep=6pt,inner sep=0pt] {
W^k={}\&\overline{\mathcal B}^k\&\BookOrthSum\&\mathsf H^k\&\BookOrthSum\&\overline{\mathcal B}_k^*\\};
\draw[decorate,decoration={brace,amplitude=4pt}] ([yshift=6pt]m-1-2.north west)--([yshift=6pt]m-1-4.north east) node[midway,above=5pt] {$\mathcal Z^k=\mathcal B_k^{*\perp}$};
\draw[decorate,decoration={brace,amplitude=4pt}] ([yshift=6pt]m-1-6.north west)--([yshift=6pt]m-1-6.north east) node[midway,above=5pt] {$\mathcal Z^{k\perp}$};
\draw[decorate,decoration={brace,mirror,amplitude=4pt}] ([yshift=-6pt]m-1-2.south west)--([yshift=-6pt]m-1-2.south east) node[midway,below=5pt] {$\mathcal Z_k^{*\perp}$};
\draw[decorate,decoration={brace,mirror,amplitude=4pt}] ([yshift=-6pt]m-1-4.south west)--([yshift=-6pt]m-1-6.south east) node[midway,below=5pt] {$\mathcal Z_k^*=\mathcal B^{k\perp}$};
\end{tikzpicture}
\end{equation}
{and}
\begin{equation}\label{eq:hodge-domain-decomposition}
V^k=\overline{\mathcal B}^k\BookOrthSum\mathsf H^k\BookOrthSum\mathcal Z^{k\perp V}.
\end{equation}
\end{Theorem}

% block: hodge-closed-complex-description page:47 type:ordinary-paragraph
{For a closed Hilbert complex we may, of course, drop the bars over the boundary spaces. In this case, the Hodge decomposition says that every element of $W$ has a unique decomposition into an element of the range of $d$, an element of the range of $d^*$, and a harmonic form, and this decomposition is orthogonal.}

% block: poincare-subheading page:47 type:subsection
\subsubsection{Poincaré {inequality}}\label{subsec:hilbert-poincare-inequality}

% block: poincare-intro page:47 type:ordinary-paragraph
{The final result we derive is the Poincaré inequality, for which it is necessary to assume that the Hilbert complex is closed.}

% block: hilbert-poincare-inequality page:47 type:theorem
\Needspace{6\baselineskip}
\begin{Theorem}{Poincaré {inequality}}\label{thm:hilbert-complex-poincare}
{Given a closed Hilbert complex, for each $k$, there exists a constant $c_P$ such that}
\begin{equation}\label{eq:hilbert-poincare-inequality}
\|z\|_V\leq c_P\|dz\|,\qquad z\in\mathcal Z^{k\perp V}.
\end{equation}
\end{Theorem}

% block: hilbert-poincare-proof page:47 type:proof
\begin{Proof}
{Since the Hilbert complex is closed, $\mathcal B^{k+1}$ is a Hilbert space. The operator $d^k$ is thus a bounded linear isomorphism between the Hilbert spaces $\mathcal Z^{k\perp V}$ and $\mathcal B^{k+1}$, and so it has a bounded inverse, which leads directly to \eqref{eq:hilbert-poincare-inequality}.}
\end{Proof}
